\chapter{Metodología ejecutada}

El proyecto se desarrolló como una investigación aplicada, cuantitativa y
experimental, acompañada por actividades de desarrollo tecnológico. Se empleó
una ruta combinada de diseño de prototipo, modelado y simulación, y preparación
de una comparación de aprendizaje por refuerzo. Cada etapa produjo artefactos
verificables: código versionado, configuraciones, pruebas, bolsas de datos,
tablas, gráficas, decisiones técnicas e informes.

\section{Correspondencia con los cuatro objetivos}

OE1 organiza la implementación de hardware y software; OE2, el modelado
cinemático y dinámico; OE3, la estrategia aprendida de coordinación; y OE4,
la evaluación mediante tres pruebas. La marcha convencional y la comparación
con y sin correcciones son procedimientos de apoyo dentro de esta estructura.

La Prueba de Margen de Estabilidad Estática debe relacionar la proyección del
centro de masa con el polígono de apoyo y conservar su evolución durante la
marcha. La Prueba de Seguimiento de Posición Articular contrasta referencias
con posiciones observadas, mediante errores por articulación. La Prueba de
Repetibilidad del Patrón de Marcha compara ensayos independientes de una misma
configuración. Hay herramientas y evidencia parcial en simulación; la aprobación
de criterios y la validación integral permanecen pendientes. Las pruebas
provocadas del supervisor verifican su reacción y no sustituyen estas tres pruebas.

\section{Plan metodológico del anteproyecto}

El anteproyecto definió una arquitectura de control de dos niveles: una marcha
nominal generada por una máquina de estados finitos, con trayectorias de los
pies y referencias articulares, y una política de aprendizaje por refuerzo que
solo produce correcciones pequeñas y acotadas sobre esas referencias, sin
controlar el PWM y subordinada a un supervisor independiente. El desarrollo se
organizó en ocho fases secuenciales con retroalimentación, de modo que cada
etapa superara sus verificaciones antes de avanzar a pruebas de mayor riesgo
para el hardware (figura~\ref{fig:metodologia_fases}).

\begin{figure}[H]
\centering
\resizebox{0.96\textwidth}{!}{%
\begin{tikzpicture}[
    node distance=0.65cm,
    block/.style={rectangle, draw, fill=blue!20, text width=10em,
        text centered, rounded corners, minimum height=3.5em},
    arrow/.style={-latex, thick}
]
\node (f1) [block] {1. Caracterización de la plataforma};
\node (f2) [block, right=of f1] {2. Modelado y simulación};
\node (f3) [block, right=of f2] {3. Electrónica y seguridad};
\node (f4) [block, right=of f3] {4. Marcha nominal};
\node (f5) [block, below=of f4] {5. Aprendizaje por refuerzo};
\node (f6) [block, left=of f5] {6. Transferencia al robot};
\node (f7) [block, left=of f6] {7. Validación comparativa};
\node (f8) [block, left=of f7] {8. Análisis y documentación};
\draw[arrow] (f1) -- (f2);
\draw[arrow] (f2) -- (f3);
\draw[arrow] (f3) -- (f4);
\draw[arrow] (f4) -- (f5);
\draw[arrow] (f5) -- (f6);
\draw[arrow] (f6) -- (f7);
\draw[arrow] (f7) -- (f8);
\draw[-latex, dashed] (f7.north) .. controls +(up:1.2cm) and +(down:1.2cm) .. (f4.south);
\draw[-latex, dashed] (f6.north) .. controls +(up:1.0cm) and +(down:1.0cm) .. (f3.south);
\end{tikzpicture}%
}
\caption{Fases de desarrollo y ciclos de retroalimentación planteados en el
anteproyecto. Fuente: elaboración propia.}
\label{fig:metodologia_fases}
\end{figure}

La tabla~\ref{tab:fases-plan-ejecucion} resume lo que cada fase preveía y cómo
se ejecutó. La transferencia progresiva prevista (sin servos, servo individual,
pata suspendida, cuatro patas suspendidas, postura con soporte, transferencia de
peso, un paso, ciclo completo y marcha continua) se cumplió hasta las cuatro
patas con el robot suspendido; los pasos con apoyo quedaron pendientes. La
validación preveía al menos diez ciclos válidos por marcha y condición, con
margen de estabilidad estática, inclinación, desplazamiento, desviación lateral,
longitud y tiempo de paso, repetibilidad y fallos; con el robot suspendido solo
fueron medibles el seguimiento articular y la repetibilidad. El anteproyecto ya
advertía que el seguimiento articular exigía medir la posición real, porque el
PCA9685 no realimenta; esa fue la razón para instrumentar las articulaciones con
potenciómetros y encoders.

\begin{table}[H]
\centering
\small
\caption{Fases del anteproyecto frente a la ejecución.}
\label{tab:fases-plan-ejecucion}
\begin{tabular}{>{\raggedright\arraybackslash}p{0.2\textwidth}>{\raggedright\arraybackslash}p{0.33\textwidth}>{\raggedright\arraybackslash}p{0.37\textwidth}}
\toprule
\textbf{Fase} & \textbf{Resultado esperado} & \textbf{Ejecución} \\
\midrule
1. Caracterización & Ficha técnica, parámetros geométricos e inerciales, inventario y riesgos & Geometría medida y aplicada al modelo; masas parciales; inercias no identificadas. \\
2. Modelado y simulación & Modelo validado y simulación de pie con referencias articulares & Cinemática, jacobianos y dinámica nominal; URDF/SDF con ROS~2 y Gazebo, MJCF en MuJoCo. \\
3. Electrónica y seguridad & Diagramas, presupuesto energético, calibración y protocolos de encendido y emergencia & Calibración articular y esquema propuesto; fuente, fusible y parada de emergencia pendientes. \\
4. Marcha nominal & Postura, paso y gateo sin aprendizaje como referencia & Gateo y paso en simulación y en el robot suspendido; trote añadido en la fase física. \\
5. Aprendizaje por refuerzo & MDP, protocolo reproducible y política correctiva validada en simulación & PPO residual con aleatorización de dominio, semillas registradas y actuador identificado (capítulo~\ref{cap:ppo-fisico}). \\
6. Transferencia & Marcha nominal y corregida ejecutadas con seguridad en el robot & Ejecutadas con el robot suspendido; topes del firmware y apagado de salidas. \\
7. Validación comparativa & Comparación estadística con y sin capa aprendida & Repetibilidad y ciclo de mejora en suspensión; menos de diez repeticiones por condición. \\
8. Análisis y documentación & Informe reproducible y documentación técnica & Este documento, repositorio y anexos de reproducción. \\
\bottomrule
\end{tabular}
\end{table}

\section{Trazabilidad y control de versiones}

El código y la documentación se mantuvieron en un repositorio Git. Para cada
ensayo se registraron la fecha, la configuración de marcha, la versión del
código, los tópicos almacenados y los hashes SHA-256 de los archivos críticos.
Las tentativas con nodos duplicados, ausencia de marcadores o pérdida de fases
se conservaron como evidencia inválida y no se mezclaron con los resultados
aceptados.

\section{Modelado y verificación de software}

Se implementaron la cinemática directa e inversa de cada pata, transformaciones
homogéneas, jacobianos y componentes de un modelo dinámico nominal. La geometría
se integró en descripciones URDF/SDF para Gazebo y MJCF para MuJoCo. La
verificación incluyó alcanzabilidad, límites articulares, continuidad de las
referencias, compilación de los paquetes ROS~2 y pruebas automatizadas.

El modelo se trató como nominal computable y no como gemelo digital, debido a
que las masas, geometría definitiva, fricción, holgura y parámetros de los
actuadores físicos todavía no se han identificado por completo.

Como preparación para esa identificación se implementó una fuente única de
perfiles que genera pares URDF/SDF--MJCF auditables. Los perfiles varían masa,
inercia, amortiguamiento, fricción articular y del suelo, holgura, retardos,
tensión y corriente. El límite del MG996R combina la envolvente de catálogo
par--velocidad con una cota de corriente; una capa ROS común aplica retardo,
zona muerta y límite de cambio de referencia en ambos simuladores. Estos
perfiles aumentan fidelidad y comparabilidad, pero sus valores continúan siendo
hipótesis hasta contrastarlos con el robot físico.

\section{Protocolo de caracterización física}

Se preparó un protocolo trazable para contrastar el modelo con el ejemplar
construido. El procedimiento separó dos actividades por seguridad. La primera,
sin energizar el robot, comprendió inventario fotográfico, inspección mecánica,
confirmación del cableado, tres repeticiones por dimensión y tres repeticiones
de cada masa disponible. Las longitudes de los eslabones se definieron entre
ejes de rotación para evitar depender de los bordes externos de las piezas.

La segunda actividad quedó condicionada a la aprobación de la alimentación y
la parada física por OE. Esta medirá individualmente el centro PWM, los límites
mecánicos conservadores, el sentido y la velocidad de cada uno de los doce
servos. No se permitirá copiar valores entre articulaciones ni habilitar el
hardware mientras existan campos sin medir.

Para una magnitud $x$ se definieron la media y la diferencia respecto del
modelo nominal $x_0$ como
\begin{equation}
  \bar{x}=\frac{x_1+x_2+x_3}{3}, \qquad \Delta x=\bar{x}-x_0.
\end{equation}
Si el intervalo entre la mayor y la menor repetición supera dos veces la
resolución del instrumento, la medición debe repetirse. El protocolo, las
fichas y los CSV se dejaron versionados antes de capturar datos. En esta etapa
se preparó el método, pero no se reportaron valores físicos como resultados.

\section{Diseño y evaluación de las marchas nominales}

Se generaron trayectorias de los pies en espacio cartesiano y se transformaron
a referencias articulares mediante la cinemática inversa propia. El gateo usó
24 referencias por ciclo, paso de 0,018~m, elevación de 0,014~m y 0,18~s por
referencia. La marcha paso usó 32 referencias, paso de 0,016~m, elevación de
0,008~m y la misma duración por referencia.

Cada ensayo comenzó con una orden explícita de marcha y terminó con una orden
\texttt{stand}. Los ciclos se delimitaron mediante las referencias y fases
registradas, no únicamente dividiendo el tiempo total por una duración nominal.
El primer ciclo se conservó como transitorio. Las métricas incluyeron avance,
velocidad, excursión lateral, altura, inclinación y continuidad articular.

\section{Contacto, IMU y estabilidad}

Se incorporaron cuatro sensores de contacto de Gazebo a 100~Hz. El monitor
consolidó el estado observado y el comparador lo contrastó con el plan publicado
por el controlador. Posteriormente se separó el estado crudo, derivado de los
mensajes y un timeout de 0,10~s, del estado filtrado. La pérdida debía persistir
0,12~s y el recontacto 0,03~s antes de modificar el estado estable.

También se integró una IMU simulada de 100~Hz. Un monitor calculó el polígono de
soporte mediante las posiciones nominales de los pies en contacto y publicó un
margen estático firmado. Debido a que el centro de masa y la geometría física
no están identificados, esta magnitud se reportó explícitamente como nominal.

\section{Seguridad}

El supervisor verificó altura, inclinación, finitud, límites y discontinuidades
de las referencias articulares. Nueve escenarios ROS~2 aislados provocaron
altura baja y alta, roll, pitch, pérdida de contacto, margen negativo, datos
vencidos, exceso de límite articular y discontinuidad. Cada criterio exigió el
motivo esperado, la activación enclavada y la orden \texttt{stand}. No se
implementó un rearme automático; contacto, margen y timeout permanecieron
deshabilitados por defecto hasta medir falsos positivos en campañas largas.

\section{Preparación experimental final}

La comparación prevista incluye, para gateo y marcha paso, cinco ensayos de 20
ciclos con control nominal y cinco con la misma marcha más corrección aprendida.
El ensayo es la unidad independiente; los ciclos no se tratarán como réplicas.
Las semillas de entrenamiento se fijaron en 11, 23, 37, 53 y 71, y los
escenarios de evaluación emparejada en 101, 202, 303, 404 y 505.

Esta comparación, prevista sobre el suelo, no se ejecutó aún. En su lugar se
realizó una comparación con el robot suspendido, descrita en el
capítulo~\ref{cap:plataforma-fisica}: varias repeticiones de 3~min por
condición, en orden alternado, con comprobación de salud tras cada caminata. Las
semillas de entrenamiento se ampliaron más allá de las cinco congeladas
(capítulo~\ref{cap:ppo-fisico}), y la selección final de las políticas se
hizo con los datos del robot.

\section{Fase experimental en el robot físico}

La fase física siguió una progresión de menor a mayor riesgo: pruebas por canal
del PCA9685; identificación de cada articulación y de sus sensores; marcha
nominal en lazo abierto con el robot suspendido; identificación del actuador con
esas caminatas; entrenamiento de políticas con el actuador identificado;
ejecución de las políticas en el robot suspendido; y, por último, un ciclo de
mejora que alterna pruebas en el robot, actualización del modelo y
reentrenamiento. En cada paso, los datos del robot (sensores y cámaras) se usaron
para detectar diferencias con la simulación antes de avanzar al siguiente.
